\section{Discussion}
\label{sec:discussion}

\subsection{What the evidence supports}
The confirmation supports a specific and bounded claim. On 31 rivers that development never used, the certificate rules were never wrong (H1), and spending bytes on witnesses of hard cases raised the share of decided hard cases over the record alone by about 0.3 at six messages per day (H2), more than the same layering on random days (H3), at a cost of 0.07 to 0.08 in the share of days whose whole record was recovered (H4). The witnesses also decided thresholds that nobody configured, which the flags decided in only 0.275 of hard cases by monotone implication (H5), and neither send-on-delta nor a tightened tolerance reached the same combination of record recovery and decisions (H6). The study does not claim that witnesses are the best way to decide thresholds that are fixed when the node is configured: flags were better there in development and in confirmation (H7).

The practical reading is a choice between three designs, depending on how fixed the thresholds are:
\begin{itemize}
  \item If the thresholds are known and stable, decision flags decide nearly every hard case (0.999 at independent loss) and recover the record almost as often as the record alone (0.877 against 0.892), but decide only 0.275 of hard cases at thresholds that were not configured, because they do not use the recovered record (the record alone decides 0.546).
  \item If thresholds may change, or if numerical extrema are wanted, hybrid 1 keeps the flag performance, adds witnesses in spare slots at no cost in record recovery or configured decisions, and raised unanticipated hard-case decisions from 0.275 to 0.629 (independent loss, descriptive), which is 0.083 above the record alone (0.097 under Gilbert loss); no design of flags plus a record certificate without witnesses was run; therefore, the share due to the hosted witnesses alone was not measured. About 78\% of development days had at least one spare slot and 21\% had four (computed at the 4.75 capacity of the flags); therefore, the benefit depends on the room that the record leaves.
  \item If the unanticipated thresholds matter most, the targeted layered design (3b rule or hybrid 2) reached 0.806, at a cost in record recovery.
\end{itemize}
This is a claim about use, not a measured benefit in the field: the unanticipated threshold sets were chosen by the analyst, and whether thresholds change in practice was not studied.

\subsection{The failed hypothesis}
H8 failed at twelve messages per day, for record only and for layered witnesses, because the upper bounds of the conditional false-certificate share exceeded their limits (0.0134 against 0.01, and 0.0215 against 0.02). The development estimates, from six sites, were lower than the cohort values by a factor of 1.2 to 15. Two lessons follow:
\begin{itemize}
  \item Corroboration: the corroboration rule compares surplus messages; therefore, its value rises with $Q$, but it cannot act when no more messages arrive than are needed, when fewer arrive and only systematic units are used, or when the record exceeds the quota and carries no parity. At $Q = 12$ the last two cases carried most of the false-certificate rate.
  \item Sample: a sample of six sites held almost none of the days with fewer messages than needed, and it understated a conditional rate that depends on the mix of days.
\end{itemize}
The consequence for use is that the CRC-16 and the free receiver rules are not shown to be sufficient at the stated limits. A longer check (CRC-24) or a message authentication code lowers the probability that a corrupted message is accepted, and it does not change the conditional share of H8. The conditional share is a property of days on which no redundancy exists to compare. A check that spans the whole record can verify a day only when every unit was decoded, which helps the days with exactly the needed number of messages; it cannot verify units that never arrived, and uncoded oversize days have no parity. For those days only more check bits per message lower the rate, and a metric that conditions on passing the check does not measure that gain; an unconditional rate at a stated bit-error rate does. None of these measures was tested.

\subsection{Limits}
\begin{itemize}
  \item Simulation only: the evidence is trace-driven replay of USGS records through modelled loss. No satellite hardware, no measured channel error statistics, and no microcontroller measurements were used.
  \item Scope: 31 rivers of the United States in one year, with only four Gulf sites; nothing is claimed for other regions, years, or deployments.
  \item Assumption R: it was off for about 15 of 38 development sites, and it was inferred from the same record that it protects. H1 is therefore a test of the code, conditional on R.
  \item Flags capacity: flags were modelled at 4.75 record values per message although the frame allows 4. This favours the flags, and it favours hybrid 1 more than the flags on the unanticipated-threshold contrast, because it enlarges the spare room for hosted witnesses; the size of the effect was not measured.
  \item Slot model: the design comparisons used an abstract slot model and slotted witnesses; hosted witnesses, the conductance pair codes, the flags, and decoding from a subset of messages had no byte-level test in the confirmation.
  \item Construction: the contrasts of H7 and the configured and record parts of H9 hold by construction; therefore, the count of supported hypotheses overstates the evidence; the informative ones are H2 to H5, the unanticipated part of H9, and H8.
  \item Integrity model: forgeries changed body bytes only; false means readings certified outside tolerance or witness codes that differ from the truth, not wrong decisions.
  \item Run record: three environment variables of the method code (the 7DADM band, the unknown-history rule, and a day limit) and the library versions were not recorded, and the freeze record was committed after the run.
  \item Analyst choices: the alternative threshold sets, the hard-case bands, the 7DADM band, and the unknown-history rule were chosen by the analyst, the last two on development data after a first run fell short.
  \item Truth: truth statements refer to the loader's record, which includes readings bridged over gaps of at most two readings, and not to the river.
  \item Witness decisions are a lower bound on what witnesses could decide, because the box rule is not exact with witnesses (Proposition 2).
  \item Integrity: the forgery model uses random content conditional on a CRC pass; an adversary needs authentication, which is out of scope.
  \item Review: the code, the proofs, and the plan were read by one separate AI session, which is an in-project reading and not an outside review, and by the project team; no outside reviewer has read them. The tests were checked by a mutation check of 18 hand-chosen faults, which found six test gaps that were then closed.
  \item Priority: the novelty statement of Section~\ref{sec:related} rests on a general web search and on sources read in part; it is not a systematic review.
\end{itemize}
